\documentclass[12pt,a4paper]{article}
\usepackage[utf8]{inputenc}
\usepackage[T1]{fontenc}
\usepackage[serbian]{babel}
\usepackage{amsmath,amssymb}
\usepackage[ruled,vlined,linesnumbered]{algorithm2e}
\usepackage{booktabs}
\usepackage{longtable}
\usepackage{geometry}
\usepackage{hyperref}
\usepackage{multirow}
\usepackage{array}
\usepackage{setspace}
\usepackage{parskip}
\usepackage{caption}

\geometry{left=2.5cm, right=2.5cm, top=2.5cm, bottom=2.5cm}
\setlength{\parindent}{0pt}
\setlength{\parskip}{6pt}

\hypersetup{colorlinks=true, linkcolor=black, urlcolor=blue, citecolor=black}

\renewcommand{\algorithmcfname}{Algoritam}

% ---------------------------------------------------------------
\title{\textbf{GRASP metaheuristika za dvostepeni transportni\\
problem sa fiksnim troškovima (TS-FCTP)}}
\author{Mateja Stojanović \\ Broj indeksa: mi251107}
\date{\today}
% ---------------------------------------------------------------

\begin{document}
\maketitle
\thispagestyle{empty}

% ---------------------------------------------------------------
\begin{abstract}
U ovom radu razmatra se Dvostepeni transportni problem sa fiksnim troškovima
(engl.\ \textit{Two-Stage Fixed-Charge Transportation Problem}, TS-FCTP),
problem kombinatorne optimizacije koji modeluje distribuciju robe kroz
dvostepeni lanac snabdevanja -- od dobavljača do distributivnih centara i od
distributivnih centara do kupaca -- uz varijabilne i fiksne troškove transporta
na svakoj vezi. Predložena je GRASP metaheuristika (engl.\ \textit{Greedy
Randomized Adaptive Search Procedure}) koja kombinuje adaptivno podešavan
parametar $\alpha$, lokalnu pretragu tipa VND sa pet okolina, egzaktno
filtriranje kandidata zasnovano na zatvorenoj formuli za trošak druge faze,
skup elitnih rešenja sa raznolikošću po Jaccardovom rastojanju, prekrajanje
putanja (engl.\ \textit{path relinking}), ILS perturbaciju elitnih rešenja
i ruin-and-recreate perturbaciju pri dugotrajnom zastoju. Implementacija je vektorizovana
(NumPy) i paralelizovana (\texttt{multiprocessing}) na nivou nezavisnih
iteracija, čime je omogućeno višestruko izvršavanje algoritma po instanci i
dobijanje statistički smislenih pokazatelja kvaliteta. Algoritam je testiran
na skupu instanci malih, srednjih i velikih dimenzija i upoređen sa rešenjima
dobijenim CPLEX solverom.
\end{abstract}
\vspace{30pt}
\textbf{Ključne reči:} dvostepeni transportni problem sa fiksnim troškovima,
GRASP, metaheuristika, egzaktno filtriranje kandidata, okolina, skup elitnih
rešenja, path relinking, perturbacija, vektorizacija, paralelizacija,
kombinatorna optimizacija.

% ===============================================================
\section{Opis problema i matematička formulacija}
% ===============================================================

\textbf{Oznake i ulazni podaci.}
Posmatra se dvostepeni lanac snabdevanja koji se sastoji od tri sloja čvorova
povezanih u dve etape transporta: $i \rightarrow j \rightarrow k$, gde je
$i \in I$ skup dobavljača (engl.\ \textit{suppliers}), $j \in J$ skup
distributivnih centara (u nastavku DC) i $k \in K$ skup kupaca (engl.\
\textit{customers}). Ulazni podaci problema su:
\begin{itemize}
  \item $s_i$ -- raspoloživa ponuda dobavljača $i \in I$;
  \item $d_k$ -- potražnja kupca $k \in K$;
  \item $b_{ij}$ -- varijabilni trošak transporta od dobavljača $i$ do DC-a $j$;
  \item $f_{ij}$ -- fiksni trošak aktiviranja veze $(i,j)$;
  \item $c_{jk}$ -- varijabilni trošak transporta od DC-a $j$ do kupca $k$;
  \item $g_{jk}$ -- fiksni trošak aktiviranja veze $(j,k)$.
\end{itemize}
Ukupna ponuda jednaka je ukupnoj potražnji, tj.\ važi
$\sum_{i \in I} s_i = \sum_{k \in K} d_k$, pa problem uvek ima dopustivo
rešenje.

\textbf{Opis problema.}
Dvostepeni transportni problem sa fiksnim troškovima (engl.\
\textit{Two-Stage Fixed-Charge Transportation Problem}, TS-FCTP) problem je
kombinatorne optimizacije koji je prvi put uveden u radu Jawahar i Balaji
\cite{jb2009}. Za svaku vezu u mreži definisana su dva troška: varijabilni,
proporcionalan količini koja se njome prevozi, i fiksni, koji se plaća samo
jednom ako se veza koristi, tj.\ ako je protok na njoj pozitivan. Cilj je
minimizacija ukupnog troška transporta uz poštovanje ograničenja kapaciteta
dobavljača i zadovoljenje potražnje svih kupaca. Odluka o tome koje veze
će biti aktivirane (koji DC-ovi i koje veze će se koristiti) ima kombinatornu
prirodu i čini problem suštinski težim od klasičnog transportnog problema, u
kome se aktiviranje veza ne naplaćuje.

\textbf{Oblasti primene.}
TS-FCTP se javlja u različitim oblastima logističkog planiranja:
\begin{itemize}
  \item \textbf{upravljanje lancima snabdevanja} -- planiranje toka robe od
    fabrika do magacina (distributivnih centara) i dalje do maloprodajnih
    objekata;
  \item \textbf{distribucija energenata} -- alokacija proizvoda od proizvodnih
    jedinica do skladišnih (distributivnih) stanica i do potrošača;
  \item \textbf{telekomunikacione mreže} -- rutiranje podataka kroz čvorove
    mreže uz fiksne troškove aktivacije veza;
  \item \textbf{humanitarna logistika} -- distribucija pomoći od centralnih
    magacina do regionalnih centara i do krajnjih korisnika.
\end{itemize}

\textbf{Promenljive.}
U matematičkoj formulaciji koriste se sledeće promenljive:
\begin{itemize}
  \item $x_{ij} \ge 0$ -- količina koja se transportuje od dobavljača $i$ do
    DC-a $j$ (tok I. etape);
  \item $y_{jk} \ge 0$ -- količina koja se transportuje od DC-a $j$ do kupca
    $k$ (tok II. etape);
  \item $z_{ij} \in \{0,1\}$ -- indikator da li se koristi veza $(i,j)$;
  \item $w_{jk} \in \{0,1\}$ -- indikator da li se koristi veza $(j,k)$.
\end{itemize}

\textbf{Matematička formulacija.}
Formulacija je preuzeta iz rada Calvete et al.\ \cite{calvete2018},
koji su, za razliku od prvobitne formulacije \cite{jb2009}, eksplicitno uveli i
ograničenje balansa protoka u distributivnim centrima. Funkcija cilja glasi:
\begin{equation} \label{eq:obj}
\sum_{i \in I}\sum_{j \in J} \left( b_{ij}x_{ij} + f_{ij}z_{ij} \right)
+ \sum_{j \in J}\sum_{k \in K} \left( c_{jk}y_{jk} + g_{jk}w_{jk} \right)
\end{equation}
Cilj problema je minimizacija funkcije cilja \eqref{eq:obj} uz ograničenja:
\begin{align}
\sum_{j \in J} x_{ij} &\le s_i, && \forall i \in I, \label{eq:sup}\\
\sum_{j \in J} y_{jk} &= d_k, && \forall k \in K, \label{eq:dem}\\
\sum_{i \in I} x_{ij} - \sum_{k \in K} y_{jk} &= 0, && \forall j \in J, \label{eq:bal}\\
x_{ij} &\le s_i\, z_{ij}, && \forall i \in I,\ \forall j \in J, \label{eq:bm1}\\
y_{jk} &\le d_k\, w_{jk}, && \forall j \in J,\ \forall k \in K, \label{eq:bm2}\\
z_{ij} &\in \{0,1\}, && \forall i \in I,\ \forall j \in J, \label{eq:domz}\\
w_{jk} &\in \{0,1\}, && \forall j \in J,\ \forall k \in K, \label{eq:domw}\\
x_{ij} &\ge 0, && \forall i \in I,\ \forall j \in J, \label{eq:domx}\\
y_{jk} &\ge 0, && \forall j \in J,\ \forall k \in K. \label{eq:domy}
\end{align}

\textbf{Značenje funkcije cilja i ograničenja.}
Funkcija cilja \eqref{eq:obj} minimizuje ukupni trošak distribucije: sabirci
$b_{ij}x_{ij}$ i $c_{jk}y_{jk}$ su varijabilni troškovi, proporcionalni
protoku, dok su $f_{ij}z_{ij}$ i $g_{jk}w_{jk}$ fiksni troškovi koji se plaćaju
samo jednom ako se odgovarajuća veza koristi. Ograničenje \eqref{eq:sup} obezbeđuje
da nijedan dobavljač ne šalje više robe nego što ima na raspolaganju.
Ograničenje \eqref{eq:dem} garantuje da se potražnja svakog kupca tačno
zadovolji. Ograničenje \eqref{eq:bal} predstavlja balans protoka u
distributivnim centrima: DC ne skladišti niti stvara robu, već je samo
prosleđuje, pa je ukupni ulaz u svaki DC jednak njegovom ukupnom izlazu.
Ograničenja \eqref{eq:bm1} i \eqref{eq:bm2} su Big-M ograničenja koja povezuju
kontinualne tokove sa binarnim indikatorima: ako je protok na vezi pozitivan,
odgovarajući indikator mora biti jednak $1$, čime se modeluje fiksni trošak
aktivacije veze.

\textbf{Složenost problema.}
Problem TS-FCTP pripada klasi NP-teških problema \cite{jb2009}: prisustvo
fiksnih troškova unosi diskontinuitete u funkciju cilja i pretvaranje
problema u problem mešovitog celobrojnog programiranja, pa egzaktno rešavanje
ne može da prati rast dimenzija realnih instanci.

\textbf{Postojeći načini rešavanja.}
U literaturi se za TS-FCTP javljaju sledeći pristupi:
\begin{itemize}
  \item \textbf{egzaktne metode} -- rešavanje MIP formulacije komercijalnim
    solverima (npr.\ IBM ILOG CPLEX) uz primenu \textit{branch-and-bound} i
    \textit{branch-and-cut} algoritama; za male instance postiže se
    optimalnost, dok se za srednje i velike instance vreme rešavanja meri
    satima \cite{calvete2018};
  \item \textbf{genetski algoritmi} -- prvi predloženi za TS-FCTP u radu
    Jawahar i Balaji \cite{jb2009}, a zatim razvijani u radu Raj i Rajendran
    \cite{raj2012};
  \item \textbf{simulirano kaljenje} -- predloženo za TS-FCTP u radu Balaji i
    Jawahar \cite{balaji2010};
  \item \textbf{hibridni genetski algoritmi i matheuristike} -- kombinuju
    evolucijsku pretragu sa egzaktnim rešavanjem potproblema ili sa snažnom
    lokalnom pretragom \cite{pop2016, calvete2018}.
\end{itemize}
U ovom radu predložena je GRASP metaheuristika koja pohlepno-nasumičnu
konstrukciju kombinuje sa promenljivom lokalnom pretragom (VND), skupom
elitnih rešenja, prekrajanjem putanja i perturbacijom; detaljan opis dat je u
narednoj sekciji.

% ===============================================================
\section{Metaheuristika GRASP i njena adaptacija}
% ===============================================================

\subsection{Osnove GRASP-a}

GRASP (\textit{Greedy Randomized Adaptive Search Procedure}) je iterativna
metaheuristika koju su uveli Feo i Resende \cite{feo1995}. Svaka iteracija se
sastoji od dve faze:
\begin{enumerate}
  \item \textbf{Faza konstrukcije}: gradi se dopustivo rešenje
    pohlepno-nasumičnim pristupom. U svakom koraku formira se
    \textit{Restricted Candidate List} (RCL) -- skup kandidata čija je ocena u
    opsegu $[\text{min},\, \text{min} + \alpha(\text{max}-\text{min})]$, gde su
    $\text{min}$ i $\text{max}$ najmanja, odnosno najveća ocena medu tekućim
    kandidatima -- i nasumično se bira jedan element RCL-a. Parametar $\alpha \in [0,1]$
    kontroliše stepen nasumičnosti: $\alpha = 0$ odgovara čistoj pohlepnoj
    heuristici, a $\alpha = 1$ nasumičnom izboru.
  \item \textbf{Faza lokalne pretrage}: dobijeno rešenje se poboljšava
    lokalnom pretragom u skupu okolnih rešenja.
\end{enumerate}
Algoritam se ponavlja unapred zadati broj iteracija i pamti globalno najbolje
pronađeno rešenje. Osnovni pseudokod GRASP-a dat je Algoritmom~\ref{alg:grasp}. Pseudokod faze
konstrukcije (GreedyRandomizedConstruction) dat je Algoritmom~\ref{alg:grc};
njegova konkretizacija za TS-FCTP data je Algoritmom~\ref{alg:construction}.

\begin{algorithm}[H]
\caption{Osnovni GRASP$(T, \alpha)$}
\label{alg:grasp}
\SetAlgoLined
$sol^* \leftarrow \varnothing$;\quad
$\text{best\_cost} \leftarrow +\infty$\;
\For{$it = 0$ \KwTo $T-1$}{
  $sol \leftarrow$ GreedyRandomizedConstruction$(\alpha)$\;
  $sol \leftarrow$ LocalSearch$(sol)$\;
  \If{$sol.\text{cost} < \text{best\_cost}$}{
    $sol^* \leftarrow sol$;\quad
    $\text{best\_cost} \leftarrow sol.\text{cost}$\;
  }
}
\Return $sol^*$
\end{algorithm}

\begin{algorithm}[H]
\caption{GreedyRandomizedConstruction$(\alpha)$}
\label{alg:grc}
\SetAlgoLined
$sol \leftarrow \varnothing$\;
\While{rešenje nije kompletno}{
  $C \leftarrow$ skup raspoloživih kandidata\;
  Izračunaj pohlepnu ocenu $g(e)$ za svako $e \in C$\;
  $g_{\min} \leftarrow \min_{e \in C} g(e)$;\quad
  $g_{\max} \leftarrow \max_{e \in C} g(e)$\;
  $\text{RCL} \leftarrow \{\, e \in C : g(e) \leq g_{\min} + \alpha(g_{\max} - g_{\min}) \,\}$\;
  Nasumično izaberi $e^*$ iz RCL\;
  Dodaj $e^*$ u $sol$ i ažuriraj skup $C$ i ocene\;
}
\Return $sol$
\end{algorithm}

\subsection{Notacija}

U nastavku su objašnjene oznake i skraćenice koje se koriste u opisu
implementacije -- GRASP pojmovi prilagođeni rešavanju problema TS-FCTP:
\begin{itemize}
  \item \textbf{DC} -- distributivni centar (čvor srednjeg sloja $J$).
  \item \textbf{I. etapa} -- transportna etapa dobavljač $\to$ DC (tokovi
    $x$, troškovi $b$, $f$).
  \item \textbf{II. etapa} -- transportna etapa DC $\to$ kupac (tokovi $y$,
    troškovi $c$, $g$).
  \item \textbf{luk} -- veza između dva čvora susednih slojeva kojom teče
    pozitivan tok, npr.\ luk $(j,k)$ znači $y_{jk} > 0$.
  \item \textbf{otvoren DC} -- DC kroz koji trenutno prolazi bar jedan
    pozitivan tok (skup \texttt{open\_dc}).
  \item \textbf{RCL} (\textit{Restricted Candidate List}) -- skup kandidata
    u fazi konstrukcije čija je ocena unutar $\alpha$-opsega od najbolje.
  \item \textbf{VND} (\textit{Variable Neighborhood Descent}) -- lokalna
    pretraga koja redom isprobava više okolina dok se ne dostigne lokalni
    optimum ni u jednoj.
  \item \textbf{first-improvement} -- pravilo prihvatanja prvog pronađenog
    pomaka koji smanjuje trošak, bez pretrage cele okoline.
  \item \textbf{skup elitnih rešenja} (\textit{elite pool}) -- kolekcija
    najboljih i međusobno raznolikih rešenja pronađenih tokom pretrage,
    izvor za prekrajanje putanja i ILS.
  \item \textbf{prekrajanje putanja} (\textit{path relinking}) -- postupak
    koji kombinuje dva elitna rešenja tako što jedno postepeno transformiše
    u drugo, razmenjujući kolone kupaca.
  \item \textbf{ILS} (\textit{Iterated Local Search}) -- perturbacija
    elitnog rešenja praćena ponovnom lokalnom pretragom.
  \item \textbf{ruin-and-recreate perturbacija} (\texttt{\_shake} u kodu) --
    perturbacija koja se aktivira pri dugotrajnom zastoju: zatvara skupo
    polovinu otvorenih DC-ova i ponovo raspoređuje njihov tok.
  \item \textbf{batch} -- grupa od $n_{\text{workers}}$ zadataka (restart,
    ILS ili prekrajanje putanja) koja se u paralelnoj verziji izvršava
    zajedno u radničkim procesima.
\end{itemize}

\subsection{Implementacija}

Implementacija (folder \texttt{GRASP\_part\_parallel/} u repozitorijumu
projekta) sastoji se iz tri fajla: \texttt{utils.py} (učitavanje instance,
infrastruktura za multi-run eksperimente i izveštaj), \texttt{GRASP.py} (ceo
algoritam) i \texttt{main.py} (CLI ulaz). U nastavku je detaljno opisana svaka
komponenta iz \texttt{GRASP.py}, a potom kompletan pseudokod algoritma.

\subsubsection{Reprezentacija rešenja -- klasa \texttt{Solution}}

Rešenje se predstavlja parom matrica toka i skupom otvorenih DC-ova:
\begin{itemize}
  \item $x \in \mathbb{Z}_{\geq 0}^{I \times J}$ -- tokovi I. etape (dobavljač $\to$ DC);
  \item $y \in \mathbb{Z}_{\geq 0}^{J \times K}$ -- tokovi II. etape (DC $\to$ kupac);
  \item $\text{open\_dc} \subseteq \{1,\ldots,J\}$ -- skup otvorenih DC-ova
    (DC $j$ je otvoren ako $\sum_i x_{ij} > 0$ ili $\sum_k y_{jk} > 0$).
\end{itemize}

Ukupni trošak rešenja (vrednost funkcije cilja) računa metoda
\texttt{compute\_cost()} kao sumu varijabilnih i fiksnih troškova u obe etape:
\begin{equation} \label{eq:cost}
\text{cost}(x,y) = \sum_{i \in I}\sum_{j \in J} b_{ij}x_{ij}
+ \sum_{j \in J}\sum_{k \in K} c_{jk}y_{jk}
+ \sum_{i \in I}\sum_{j \in J} f_{ij}\,[x_{ij} > 0]
+ \sum_{j \in J}\sum_{k \in K} g_{jk}\,[y_{jk} > 0],
\end{equation}
gde je $[P]$ Iversonov simbol: $[P] = 1$ ako je tvrđenje $P$ tačno, a $0$ u
suprotnom. Fiksni trošak veze se dakle plaća tačno kada je protok na njoj
pozitivan. Svaki poziv \texttt{eval\_cost()} -- stvarno izračunavanje
vrednosti funkcije cilja -- uvećava globalni brojač $eval$ (broj poziva
funkcije cilja). Metoda \texttt{compute\_cost()} je omotač koji poziva
\texttt{eval\_cost()} i zatim osvežava \texttt{open\_dc}; u $eval$ se ne beleži
broj poziva \texttt{compute\_cost()}, jer lokalna pretraga u prihvatnoj probi
poziva \texttt{eval\_cost()} i bez osvežavanja \texttt{open\_dc}. Metoda
\texttt{is\_valid()} proverava dopustivost rešenja: potražnja svakog kupca,
kapacitet svakog dobavljača i balans protoka u svakom DC-u.

Ključna arhitekturalna odluka: tokom pretrage se direktno modifikuju
isključivo $y$-tokovi (II. etapa); $x$-tokovi (I. etapa) se pri svakoj promeni
$y$-toka iznova rekonstruišu pohlepnom procedurom \textsc{Repair}. Iako se
potezi u lokalnoj pretrazi zovu \texttt{\_add\_dc}, \texttt{\_remove\_dc} i
\texttt{\_swap\_dc}, oni ne otvaraju/zatvaraju DC-ove nezavisno od tokova: DC
je otvoren ako i samo ako kroz njega teče pozitivan tok, pa se svaki od tih
poteza u stvari svodi na premeštanje $y$-tokova kojima se pojedini DC-ovi
pune, odnosno prazne. Ovaj pristup smanjuje dimenzionalnost prostora pretrage
jer algoritam nikad eksplicitno ne upravlja $x$-promenljivama.

\subsubsection{Procedura \textsc{Repair}}

Za svaki DC $j$ sa pozitivnim ukupnim prilivom ($\text{demand}_j =
\sum_k y_{jk}$), dobavljači se rangiraju po celobrojnoj oceni koja kombinuje
varijabilni i amortizovani fiksni trošak,
\begin{equation} \label{eq:score}
\text{score}_i = b_{ij}\cdot\text{demand}_j + f_{ij},
\end{equation}
a zatim se tokovi pohlepno dodeljuju u tom redosledu dok se zahtev DC-a ne
zadovolji (pri izjednačenim ocenama prednost ima dobavljač sa najvećim
preostalim kapacitetom). Ovo je celobrojna varijanta ocene
$b_{ij} + f_{ij}/\text{demand}_j$ -- pošto je $\text{demand}_j > 0$ fiksno za
dati DC, isto rangiranje se dobija skaliranjem sa $\text{demand}_j$, čime se
izbegava deljenje u pokretnom zarezu i gubitak preciznosti. Ako neki DC ne
može biti u potpunosti snabdeven (nedostatak ukupnog kapaciteta), procedura
vraća neuspeh. Pseudokod je dat Algoritmom~\ref{alg:repair}.

\begin{algorithm}[H]
\caption{Repair$(sol, D)$}
\label{alg:repair}
\SetAlgoLined
$sol.x \leftarrow \mathbf{0}$\;
$\text{rem}\_s \leftarrow D.s.copy()$\;
\For{$j = 1$ \KwTo $J$}{
  $\text{demand}_j \leftarrow \sum_k y_{jk}$\;
  \lIf{$\text{demand}_j = 0$}{\textbf{continue}}
  Sortiraj dobavljače po $\text{score}_i = b_{ij}\cdot\text{demand}_j + f_{ij}$\;
  $\text{rem} \leftarrow \text{demand}_j$\;
  \For{svakog $i$ u sortiranom redosledu}{
    $q \leftarrow \min(\text{rem}\_s[i],\, \text{rem})$\;
    $x_{ij} \mathrel{+}= q$;\quad
    $\text{rem}\_s[i] \mathrel{-}= q$;\quad
    $\text{rem} \mathrel{-}= q$\;
    \lIf{$\text{rem} = 0$}{\textbf{break}}
  }
  \lIf{$\text{rem} > 0$}{\Return \textit{neuspešno}}
}
\Return \textit{uspešno}
\end{algorithm}

\subsubsection{Faza konstrukcije -- klasa \texttt{Construction}}

Kupci se obrađuju u opadajućem redosledu po potražnji. Za svakog kupca $k$, u
svakom koraku se formira RCL od parova $(i,j)$ sa ocenom
\begin{equation} \label{eq:score_cons}
\text{score}(i,j) = b_{ij} + c_{jk}
  + \frac{f_{ij}}{need}\cdot[x_{ij} = 0]
  + \frac{g_{jk}}{need}\cdot[y_{jk} = 0],
\end{equation}
gde je $need$ preostala potražnja kupca $k$ u tekućem koraku, a
$[x_{ij} = 0]$ i $[y_{jk} = 0]$ indikatori koji dodaju amortizovani fiksni
trošak samo za veze koje još nisu aktivirane. RCL sadrži sve parove čija je
ocena $\leq \text{score}_{\min} + \alpha(\text{score}_{\max} -
\text{score}_{\min})$, a iz njega se nasumično bira par kome se dodeljuje tok
$q = \min(\text{rem\_s}_i, need)$.

Diversifikacija na nivou skupa otvorenih DC-ova postiže se parametrom
\texttt{random\_dc\_mask}: sa verovatnoćom $p = 0{,}5$, konstrukcija (i naredna
lokalna pretraga) ograničava se na nasumično izabran podskup DC-ova čija je
veličina između $\lfloor 0{,}35\,J \rfloor$ (najmanje $1$) i $J-1$. Ovim se
razbija determinizam pohlepne
konstrukcije -- bez maske, jeftiniji DC-ovi pobeđuju pri svakom restartu i
sva rešenja konvergiraju u isti skup otvorenih DC-ova. Pseudokod je dat
Algoritmom~\ref{alg:construction}.

\begin{algorithm}[H]
\caption{GreedyRandomizedConstruction$(D, \alpha, \text{allowed\_dc})$}
\label{alg:construction}
\SetAlgoLined
$sol \leftarrow$ prazno rešenje\;
\For{$k$ u opadajućem redosledu po $d_k$}{
  $need \leftarrow d_k$\;
  \While{$need > 0$}{
    Vektorizovano izračunaj $I \times J$ matricu ocena $\text{score}(i,j)$\;
    Maskiraj dobavljače sa $\text{rem}\_s[i] = 0$ i DC-ove van
    $\text{allowed\_dc}$ (ocena $= \infty$)\;
    $\text{score}_{\min} \leftarrow \min$, $\text{score}_{\max} \leftarrow \max$ (nad konačnim vrednostima)\;
    $\text{RCL} \leftarrow \{(i,j): \text{score}(i,j) \leq \text{score}_{\min} + \alpha(\text{score}_{\max}-\text{score}_{\min})\}$\;
    Nasumično izaberi $(i^*, j^*)$ iz RCL\;
    $q \leftarrow \min(\text{rem}\_s[i^*],\, need)$\;
    $x_{i^*j^*} \mathrel{+}= q$;\quad
    $y_{j^*k} \mathrel{+}= q$;\quad
    $\text{rem}\_s[i^*] \mathrel{-}= q$;\quad
    $need \mathrel{-}= q$\;
  }
}
Izračunaj cenu $sol$\;
\Return $sol$
\end{algorithm}

\subsubsection{Lokalna pretraga -- klasa \texttt{LocalSearch} (VND)}

Primenjuje se \textit{Variable Neighborhood Descent} (VND) sa pet okolina,
metodom \textit{first improvement} -- prvi pronađeni pomak koji smanjuje cenu
prihvata se i pretraga se vraća na prvu okolinu. Ako nijedna okolina ne
poboljša rešenje, brojač $no\_improve$ se uvećava; posle
$\text{max\_no\_improve} = 10$ neuspešnih ciklusa pretraga se zaustavlja.
Redosled okolina nije slučajan: okoline su poredane od manje ka sveobuhvatnijoj
izmeni rešenja. \texttt{\_reassign\_customer} ne menja skup otvorenih DC-ova
(kupac se seli na već otvoren DC); \texttt{\_add\_dc}, \texttt{\_remove\_dc} i
\texttt{\_swap\_dc} menjaju sam skup otvorenih DC-ova, koji nosi ključnu
kombinatornu odluku problema; \texttt{\_shift\_flow} na kraju fino podešava
raspodelu cepanjem kupca i isplati se jedino kroz I. etapu. Kako VND posle
svakog poboljšanja vraća na prvu okolinu, manje invazivni (i jeftiniji) potezi
isprobavaju se češće, a sveobuhvatniji tek kada jednostavniji ne donose
napredak.
\begin{enumerate}
  \item \textbf{Preraspoređivanje kupca} (\texttt{\_reassign\_customer}): ovaj
    potez pokušava da kupca $k$, koga trenutno može opsluživati više DC-ova,
    prebaci u celosti na jedan DC $j_{\text{to}}$. Prvo se DC-ovi poređaju po
    tome koliko bi koštalo opsluživanje tog kupca sa svakog od njih (trošak
    II. etape); zatim se od najviše $8$ njih (najboljih $6$ plus $2$
    nasumična, radi istraživanja) redom proba svaki, i prihvata se prvi koji
    stvarno smanji ukupnu cenu rešenja. Rangiranje po trošku II. etape samo
    bira redosled isprobavanja -- ne odbacuje kandidate unapred, jer ušteda
    ponekad dolazi iz I. etape, a ne iz II.
  \item \textbf{Dodavanje DC-a} (\texttt{\_add\_dc}): ovaj potez pokušava da
    otvori jedan zatvoreni DC; od svih trenutno zatvorenih DC-ova isprobava se
    najviše $36$ kao kandidati. Za razliku od prethodnog poteza, ovde se ne
    premešta ceo kupac odjednom, već pojedinačne veze $(j,k)$ -- tako kupac
    koji se već deli između više DC-ova zadržava ostatak svoje podele, a menja
    se samo jedna njegova veza. Za svaki kandidatski DC $j_{\text{new}}$ prvo
    se za svaku postojeću vezu izračuna koliko bi (u trošku II. etape) koštalo
    da se ta veza umesto na svoj trenutni DC preusmeri na $j_{\text{new}}$, a
    zatim se probaju dve varijante otvaranja tog DC-a (prva koja da jeftinije
    rešenje se prihvata):
    \begin{enumerate}
      \item[(i)] preusmere se \emph{sve} veze kod kojih je to preusmeravanje
        samo po sebi jeftinije u II. etapi;
      \item[(ii)] preusmere se samo $3$ veze kod kojih je to preusmeravanje
        najmanje skuplje (bez obzira da li je uopšte jeftinije) -- ova
        eksplorativna varijanta otvara DC i kada II. etapa sama po sebi ne
        opravdava potez, jer se ušteda može javiti tek kroz I. etapu.
    \end{enumerate}
  \item \textbf{Uklanjanje DC-a} (\texttt{\_remove\_dc}): isprobava se do
    $\text{REMOVE\_MAX\_CANDIDATES} = 20$ otvorenih DC-ova za zatvaranje; za
    svakog pogođenog kupca vektorizovano se bira najjeftinija preostala
    alternativa.
  \item \textbf{Zamena DC-a} (\texttt{\_swap\_dc}): jedan otvoreni DC se zamenjuje
    jednim zatvorenim; oslobođeni kupci se preraspodeljuju na novootvoreni ili
    na najjeftiniji od preostalih otvorenih DC-ova (videti i
    odeljak~\ref{sec:egzaktno-filtriranje}).
  \item \textbf{Pomeranje dela toka} (\texttt{\_shift\_flow}): dok prethodni
    potezi uvek premeštaju \emph{celu} vezu ili celog kupca, ovaj potez sme da
    premesti samo \emph{deo} količine sa postojeće veze između DC-a $j$ i
    kupca $k$ na drugi, ciljni DC $j_{\text{to}}$, ostavljajući ostatak na
    $j$ -- kupac $k$ posle toga ima dve veze umesto jedne, tj.\ postaje
    podeljen. To je jedina okolina koja
    ovako deli kupca; sve ostale ili premeštaju čitavu vezu ili je ostavljaju
    netaknutom. Za količinu koja se premešta isprobavaju se četiri opcije:
    cela veza, polovina veze, ili jedna od dve količine izračunate tako da
    posle premeštanja protok kroz izvorni odnosno ciljni DC postane
    umnožak kapaciteta jednog dobavljača -- to poslednje je bitno jer
    \textsc{Repair} snabdeva DC redom, od najjeftinijeg dobavljača, uzimajući
    cele kapacitete dobavljača (svaki iznosi 80). DC pritom zaista može imati
    više dobavljača, ali svaki je iskorišćen u celosti. Ako protok DC-a nije
    umnožak broja 80, na kraju ostaje ostatak (manji od 80) koji mora da
    pokrije još jedan, najskuplji dobavljač -- angažovan samo za taj ostatak,
    uz svoj fiksni trošak. Kad je protok umnožak broja 80, ostatka nema, pa
    taj dodatni dobavljač za taj DC nije potreban. Potez uzorkuje najviše $200$ veza,
    za svaku računa tačnu promenu troška II. etape za sve četiri opcije, i
    redom (od najjeftinije ka najskupljoj) proba par (veza, količina, ciljni
    DC) dok neki ne smanji ukupnu cenu. Deljenje kupca samo po sebi nikad ne
    smanjuje trošak II. etape (koji je linearan po količini plus fiksna
    doplata), pa se ovaj potez isplati jedino kroz I. etapu -- svaki
    dobavljač ima ograničen kapacitet, pa kada se njegovi najjeftiniji
    dobavljači potroše, dalje snabdevanje jednog DC-a postaje sve skuplje;
    deljenje kupca na dva DC-a može tada biti jeftinije nego snabdevati ga
    u celosti sa jednog. Pomeranje toka u I. etapi ne definiše se kao zaseban
    potez: $x$-tokovi se ne menjaju direktno, već se posle svake promene $y$
    iznova računaju u \textsc{Repair}-u (najjeftiniji dobavljač prvi), pa svaka
    ručno rađena izmena $x$ bez promene $y$ ne bi preživela naredni poziv
    \textsc{Repair}-a. Čisto poboljšanje I. etape uz fiksirano $y$ već
    sistematski radi \texttt{polish\_x} na kraju VND ciklusa, pa bi dodatni
    potez između dva dobavljača bio suvišan.
\end{enumerate}
Razlog za dodatni korak je sama procedura \textsc{Repair}: ona obrađuje
DC-ove redom po indeksu, pa DC $1$ uvek prvi bira najjeftinije dobavljače, a
DC-ovi koji dođu na red kasnije nasleđuju ono što je preostalo -- rezultat
nije nužno ni lokalno optimalan u I. etapi. Na konvergenciji VND ciklusa
zato se poziva procedura \texttt{polish\_x}, koja uz fiksirano $y$ dodatno
poboljšava trošak I. etape. Ona radi isključivo u matrici $x$: bira dva
postojeća luka $(i_1,j_1)$ i $(i_2,j_2)$ sa različitim dobavljačima i
različitim DC-ovima, i razmenjuje po njima jednaku količinu $d$ unakrsno --
$x_{i_1j_1}$ i $x_{i_2j_2}$ se umanje za $d$, a $x_{i_1j_2}$ i $x_{i_2j_1}$
uvećaju za $d$ (tzv.\ četvorougaona razmena). Količina $d$ bira se kao
$d = \min(x_{i_1j_1},\, x_{i_2j_2})$ -- najveća vrednost koja se može
oduzeti od oba luka a da nijedan tok ne postane negativan; luk s manjim
tokom pritom se u potpunosti gasi, čime se štedi i njegov fiksni trošak.
Ovakva razmena ne menja ni
koliko svaki dobavljač ukupno isporučuje ni koliko svaki DC ukupno prima, pa
$y$ i dopustivost rešenja ostaju netaknuti -- menja se samo \emph{koji}
dobavljač snabdeva koji DC. Za svaki par lukova unapred se, bez pravljenja
probnog rešenja, izračuna tačna promena troška takve razmene, i uvek se bira
par sa najvećom uštedom (best-improvement), sve dok postoji razmena koja
smanjuje trošak ili dok se ne dostigne najviše $200$ pokušaja. Pošto broj
korišćenih lukova raste sa veličinom instance, kandidati se uzorkuju ako ih
ima više od $\text{max\_arcs} = \max(600,\, 2(I+J))$. Pseudokod je dat
Algoritmom~\ref{alg:vnd}.

\begin{algorithm}[H]
\caption{VND$(D,\, sol,\, \text{allowed\_dc},\, \text{max\_no\_improve}=10)$}
\label{alg:vnd}
\SetAlgoLined
\textsc{Repair}$(sol, D)$;\quad Izračunaj cenu $sol$\;
$\text{no\_improve} \leftarrow 0$\;
\While{$\text{no\_improve} < \text{max\_no\_improve}$}{
  \uIf{ReassignCustomer$(sol)$}{$\text{no\_improve} \leftarrow 0$; \textbf{continue}}
  \uElseIf{AddDC$(sol)$}{$\text{no\_improve} \leftarrow 0$; \textbf{continue}}
  \uElseIf{RemoveDC$(sol)$}{$\text{no\_improve} \leftarrow 0$; \textbf{continue}}
  \uElseIf{SwapDC$(sol)$}{$\text{no\_improve} \leftarrow 0$; \textbf{continue}}
  \uElseIf{ShiftFlow$(sol)$}{$\text{no\_improve} \leftarrow 0$; \textbf{continue}}
  \Else{$\text{no\_improve} \mathrel{+}= 1$}
}
$\text{polish\_x}(sol, D)$\;
\Return $sol$
\end{algorithm}

\subsubsection{Skup elitnih rešenja -- klasa \texttt{ElitePool}}

Čuva se skup od do $12$ najkvalitetnijih i međusobno raznolikih rešenja.
Raznolikost se meri Jaccardovim rastojanjem između skupova otvorenih DC-ova
dva rešenja $a$ i $b$:
\begin{equation} \label{eq:jaccard}
d(a,b) = 1 - \frac{|\,a.\text{open\_dc} \cap b.\text{open\_dc}\,|}
               {|\,a.\text{open\_dc} \cup b.\text{open\_dc}\,|},
\end{equation}
sa pragom $\text{min\_dist} = 0{,}2$. Pravilo ažuriranja prati standardnu
politiku elitnog skupa: dok skup nije pun, odbijaju se samo rešenja sa istim
skupom otvorenih DC-ova (ukoliko nisu jeftinija od tog člana);
kada je pun, kandidat koji je suviše blizu nekom članu (\texttt{open\_dc}
rastojanje $\leq \text{min\_dist}$) takmiči se samo sa \textit{tim} najbližim
članom i može ga zameniti jedino ako ima nižu cenu; kandidat dovoljno udaljen
od svih članova zamenjuje najskupljeg člana ako je jeftiniji od njega. Time
jeftino rešenje nikad ne može da istisne raznolikost skupa -- može da zameni
samo sopstvenog najbližeg suseda. Metode \texttt{best()} i \texttt{random()}
vraćaju najjeftinije, odnosno nasumično rešenje iz skupa, i koriste se kao ulaz
za prekrajanje putanja.

\subsubsection{Prekrajanje putanja -- klasa \texttt{PathRelinking}}

Svakih $8$ iteracija uzimaju se dva rešenja iz skupa elitnih rešenja -- zovimo
ih $a$ (najbolje) i $b$ (nasumično izabrano) -- a između njih se izvodi
\textit{prekrajanje putanja} razmenom kolona kupaca. Ključno zapažanje: u
svakom dopustivom rešenju kolona $y[:,k]$ sabira tačno $d_k$, pa zamena
kolone jednog kupca u rešenju $a$ kolonom tog istog kupca iz rešenja $b$
zadržava II. etapu tačno dopustivom -- treba samo ponovo izgraditi I. etapu
procedurom \textsc{Repair}.

Postupak polazi od $a$ i, kolonu po kolonu, pretvara ga u $b$: neka je $cur$
tekuće (menjano) rešenje, koje na početku jednako $a$, a $b$ nepromenjeni
cilj kome $cur$ teži. U svakom koraku se od kupaca čije se kolone u $cur$ i
$b$ još razlikuju nasumično izabere do $\text{max\_cand} = 8$ njih; za
svakog od njih se probno napravi kopija $cur$ u kojoj je njegova kolona
zamenjena odgovarajućom kolonom iz $b$, i izračuna cena te kopije posle
\textsc{Repair}-a. Od tih probnih zamena, u $cur$ se stvarno primeni ona sa
najnižom cenom -- bez obzira da li je ta cena niža ili viša od cene
trenutnog $cur$ pre zamene. Postupak, drugim rečima, ne prihvata samo
poboljšanja: on nastavlja i kad ga sledeći korak privremeno poskupi, jer je
cilj da $cur$ stigne do $b$, a ne da se zaustavi na prvom lokalnom
minimumu. Zato se posebno pamti \emph{najjeftinije} rešenje viđeno
duž celog postupka -- ono se na kraju vraća kao rezultat, bez obzira gde se
postupak zaustavio. Postupak se zaustavlja posle najviše
$\text{max\_steps} = 15$ koraka ili kada $cur$ postane jednako $b$ (sve
kolone usaglašene).

Ovo se radi u \textit{oba} smera -- jednom kao $a \to b$ ($cur$ kreće od $a$
i cilja $b$), jednom kao $b \to a$ ($cur$ kreće od $b$ i cilja $a$) -- pošto
su međurešenja na ta dva puta različita; čuva se najjeftinija tačka viđena
na bilo kojem od njih, i ona se na kraju finalizuje lokalnom pretragom
(VND). Skupovi otvorenih DC-ova ne igraju ulogu u postupku --
unija i presek otvorenih DC-ova roditelja ne koriste se, već se putanja definiše
isključivo razlikama u kolonama kupaca. Pseudokod je dat
Algoritmom~\ref{alg:pr}.

\begin{algorithm}[H]
\caption{PathRelinking$(a, b, D)$}
\label{alg:pr}
\SetAlgoLined
$\text{best} \leftarrow \varnothing$;\quad
$\text{best\_cost} \leftarrow +\infty$\;
\For{$(src, dst) \in \{(a,b),\, (b,a)\}$}{
  $cur \leftarrow src.\text{copy}()$\;
  $diff \leftarrow \{ k : cur.y[:,k] \neq dst.y[:,k] \}$\;
  $\text{steps} \leftarrow 0$\;
  \While{$diff \neq \varnothing$ \textbf{i} $\text{steps} < \text{max\_steps}$}{
    $\text{cands} \leftarrow$ nasumični podskup od $\le \text{max\_cand}$ kupaca iz $diff$\;
    $k^* \leftarrow \arg\min_{k \in \text{cands}}\; \text{cost}\big(\text{Repair}(cur,\,D)
    \text{ posle } cur.y[:,k] \leftarrow dst.y[:,k]\big)$\;
    $cur.y[:,k^*] \leftarrow dst.y[:,k^*]$\;
    \textsc{Repair}$(cur, D)$;\quad izračunaj cenu $cur$\;
    $diff \leftarrow diff \setminus \{k^*\}$;\quad
    $\text{steps} \mathrel{+}= 1$\;
    \If{$cur.\text{cost} < \text{best\_cost}$}{
      $\text{best} \leftarrow cur.\text{copy}()$;\quad
      $\text{best\_cost} \leftarrow cur.\text{cost}$\;
    }
  }
}
\Return VND$(D, \text{best})$
\end{algorithm}

U paralelnoj verziji prekrajanje putanja se izvršava u radnicima, a ne u glavnom
procesu (videti odeljak~\ref{sec:vektorizacija}): merenja su pokazala da iznosi oko $30\%$
ukupnog vremena izvršavanja, pa bi sekvencijalno izvršavanje držalo sve radnike
neaktivnima dok glavni proces ne završi.

\subsubsection{Perturbacija i intenzifikacija}

Perturbacija se koristi na dva nezavisna nivoa: ILS perturbacija je
\textit{intenzifikacioni} mehanizam (ponovni spust elitnog rešenja u njegovom
DC-okruženju), dok je ruin-and-recreate perturbacija \textit{diversifikacioni}
mehanizam koji se aktivira tek pri dugotrajnom zastoju.
\begin{itemize}
  \item \textbf{ILS perturbacija}: sa verovatnoćom $p_{seed} = 0{,}5$ iteracija
    ne kreće iz praznog rešenja, već se nasumično odabrano elitno rešenje
    (\textit{seed}) perturbuje -- nasumičan broj otvorenih DC-ova (do otprilike
    trećine) se zatvori, a njihovi lukovi se \textit{nasumično} preliju na druge
    DC-ove (namerno ne pohlepno, jer bi pohlepno presmeravanje bilo samo još
    jedan lokalni potez koji bi VND odmah poništio) -- a zatim se spušta
    lokalnom pretragom. Ponovni spust se ograničava maskom
    \texttt{elite\_dc\_mask}, koja lokalnu pretragu zadržava na skupu otvorenih
    DC-ova semena uvećanom za $n_{extra} = 1$ nasumično izabrani zatvoreni DC.
    Bez maske bi se svaka perturbacija elitnog rešenja vratila u isti generički
    skup otvorenih DC-ova, pa skup DC-ova koji prekrajanje putanja pokušava da
    pronađe nikad ne bi bio intenzifikovan.
  \item \textbf{Ruin-and-recreate perturbacija}: ako se $20$ uzastopnih
    iteracija ne pronađe poboljšanje najboljeg rešenja, DC-ovi se rangiraju po
    prosečnom varijabilnom trošku po jedinici protoka, uklanja se $50\%$
    najskupljih, a kupci koji su im bili dodeljeni presmeravaju se na preostale
    DC-ove po najnižoj oceni troška. Posle \textsc{Repair}-a i provere
    dopustivosti primenjuje se VND; poboljšanje se prihvata kao novo najbolje,
    a i samo rešenje se, kao potpuno spušteno, dodaje u skup elitnih rešenja
    bez obzira na poboljšanje. Brojač zastoja se resetuje posle svakog pokušaja
    (bez obzira na uspeh), čime se sprečavaju uzastopne perturbacije iz istog
    stanja.
\end{itemize}

\subsubsection{Adaptivno podešavanje parametra $\alpha$}

Ako je prošlo više od $15$ iteracija bez poboljšanja najboljeg rešenja,
$\alpha$ se uzorkuje iz $\mathcal{U}(0{,}4,\, 0{,}7)$ radi diversifikacije
pretrage (veći RCL, više nasumičnosti). Inače, $\alpha \in
\mathcal{U}(0{,}05,\, 0{,}5)$ radi intenzifikacije (manji RCL, bliže pohlepnom
izboru). Reč je o \emph{reaktivnom podešavanju zasnovanom na stagnaciji}, a ne
o klasičnom Reactive GRASP-u, u kome se verovatnoće pojedinih vrednosti
$\alpha$ prilagođavaju prema kvalitetu rešenja koja su proizvele; ovde se
$\alpha$ bira iz dva fiksna opsega, u zavisnosti od toga da li pretraga
napreduje.

\subsubsection{Egzaktno filtriranje kandidata}
\label{sec:egzaktno-filtriranje}

Puna evaluacija jednog kandidata je skupa: treba napraviti probno rešenje,
pozvati \textsc{Repair} (koji je $O(I \times J)$) i tek onda izračunati cenu.
Kad okolina ima na desetine ili stotine kandidata po pozivu, provera svih na
ovaj način je najskuplji deo pretrage. Zato se u više okolina
(\texttt{\_swap\_dc}, \texttt{\_shift\_flow}, \texttt{\_reassign\_customer})
koristi \textbf{dvostepena evaluacija}: prvo se svi kandidati jeftino
poređaju po \emph{tačnoj} promeni troška II. etape -- koja se, za razliku od
troška I. etape, može izračunati direktno iz formule, bez pravljenja probnog
rešenja i bez \textsc{Repair}-a -- a zatim se \textsc{Repair} i puna cena
računaju samo za mali broj najbolje rangiranih kandidata, plus nekoliko
nasumičnih. Ovi nasumični kandidati su tu jer rangiranje vidi samo II. etapu:
kandidat koji izgleda loše u II. etapi može ipak biti dobar zbog uštede u I.
etapi (npr.\ rasterećenja skupog dobavljača), pa rangiranje sme da bira
\emph{redosled} provere, ali ne sme unapred da odbaci kandidate.

Ilustracije radi, posmatrajmo kako to radi okolina zamene DC-a
(\texttt{\_swap\_dc}): jedan otvoreni DC $j'$ se zatvara, a jedan zatvoreni
$j''$ otvara, pa svi kupci koje je opsluživao $j'$ moraju da se presele na
neki drugi otvoreni DC ili na $j''$. Za dati par $(j', j'')$, trošak II.
etape te zamene se može izračunati unapred, bez ikakve simulacije: za
svakog kupca $k$ koga je opsluživao $j'$ dovoljno je znati koliko bi koštalo
da ga umesto toga opslužuje bilo koji drugi trenutno otvoren DC ili $j''$
(trošak $= $ varijabilni trošak $c_{jk}$ puta količina, plus fiksni trošak
$g_{jk}$ ako ta veza još nije u upotrebi), i uzeti najjeftiniju opciju za tog
kupca. Zbir tih najjeftinijih opcija po svim oslobođenim kupcima je egzaktan
trošak II. etape cele zamene -- razlika ($\Delta$) tog troška u odnosu na
sadašnji trošak istih kupaca je ključ po kome se par $(j', j'')$ rangira.
Formalno, uz $q_{j'k} = y_{j'k}$ (koliko kupac $k$ trenutno dobija od $j'$) i
$K_{j'}$ (skup kupaca koje opslužuje $j'$), trošak II. etape opsluživanja
kupca $k$ iz DC-a $j$ je
\begin{equation} \label{eq:stage2}
s_2(j,k) = c_{jk}\, q_{j'k} + g_{jk}\,[y_{jk} = 0],
\end{equation}
gde indikator $[y_{jk} = 0]$ dodaje fiksni trošak samo ako ciljni luk nije
već u upotrebi, pa je egzaktni trošak II. etape cele zamene
\begin{equation} \label{eq:swapdelta}
C^{(2)}_{j',j''} = \sum_{k \in K_{j'}}
\min\Big( s_2(j'',k),\ \min_{j \in J_{\text{open}} \setminus \{j'\}} s_2(j,k) \Big).
\end{equation}
Svi parovi (otvoren, zatvoren) rangiraju se po $\Delta = C^{(2)}_{j',j''} -
C^{(2)}_{j'}$, gde je $C^{(2)}_{j'}$ tekući trošak istih kupaca dok ih još
opslužuje $j'$. Ovo se ne mora ponavljati za svaki par posebno: oslobođeni
kupci zavise samo od $j'$, pa se izračunavanje grupiše po $j'$, i cela
$|J_{\text{open}}| \times |J_{\text{closed}}|$ mreža parova košta samo
$|J_{\text{open}}|$ vektorizovanih prolaza. U punoj evaluaciji (probno
rešenje, \textsc{Repair}, prava cena) izvršavaju se samo najbolje rangiranih
$\text{SWAP\_MAX\_TRIALS} = 12$ parova, plus $\text{SWAP\_EXPLORE} = 2$
nasumično izabrana iz ostatka od $\text{SWAP\_MAX\_PAIRS} = 150$ rangiranih.

Isto načelo primenjuje se i u ostalim okolinama, sa budžetima definisanim kao
atributi klase (npr.\ $\text{SHIFT\_MAX\_ARCS} = 200$,
$\text{SHIFT\_MAX\_TRIALS} = 40$ za okolinu pomeranja dela toka). Merenja
pokazuju da filtriranje ne pogoršava kvalitet rešenja, a ubrzava lokalnu
pretragu za $10$--$20\%$ na instancama gde je budžet aktivan.

\subsubsection{Vektorizacija i paralelizacija}
\label{sec:vektorizacija}

Svi najzahtevniji delovi algoritma -- izračunavanje matrica ocena u konstrukciji,
evaluacija kandidata u okolinama i \texttt{compute\_cost()} -- vektorizovani
su preko NumPy broadcast-a i boolean maski, bez promene rezultata pretrage
(ista RCL/$\alpha$ logika i ista \textit{first-improvement} pravila, samo
izračunata bržim putem).

GRASP petlja je inherentno sekvencijalna (skup elitnih rešenja, praćenje
najboljeg rešenja, prekrajanje putanja i perturbacija zavise od deljenog stanja),
ali konstrukcija i lokalna pretraga jedne iteracije zavise samo od parametra
$\alpha$ (odnosno odabranog semena za ILS granu) i od nepromenjenih podataka
instance. Zato se batch od $n_{\text{workers}}$ nezavisnih zadataka izvršava
paralelno preko \texttt{multiprocessing.Pool}:
\begin{itemize}
  \item Podaci instance šalju se svakom radniku tačno jednom, preko
    \texttt{initializer} funkcije; radnici vraćaju samo rezultat ($x$, $y$,
    cena, \texttt{open\_dc}, broj poziva funkcije cilja).
  \item Svaki zadatak je \textit{tagovana} n-torka: \texttt{('ls', $\alpha$,
    None)} -- običan GRASP restart (konstrukcija + VND pod nasumičnom maskom
    DC-ova); \texttt{('ls', $\alpha$, seed\_y)} -- ILS iteracija (perturbacija
    elitnog rešenja i ponovni spust pod \texttt{elite\_dc\_mask}); ili
    \texttt{('pr', ...)} -- prekrajanje putanja između dva elitna rešenja.
  \item Sekvencijalni deo -- ažuriranje skupa elitnih rešenja, praćenje
    najboljeg rešenja, prekrajanje putanja i perturbacija -- primenjuje se na
    rezultate iz batch-a istim redosledom kao u serijskoj verziji. To nije
    bit-identično serijskoj petlji: parametar $\alpha$, ILS semena i roditelji
    za prekrajanje putanja fiksirani su na \textit{početku} batch-a, pa
    konstrukcija$+$VND (odnosno prekrajanje putanja) rade nad stanjem skupa
    elitnih rešenja kakvo je bilo pre nego što je batch započeo, a ne nad
    stanjem koje bi u serijskom režimu već apsorbovalo ranije iteracije
    batch-a.
  \item Ako je $n_{\text{workers}} \le 1$, koristi se u potpunosti originalna
    serijska petlja -- ponašanje je tada identično algoritmu bez
    paralelizacije.
\end{itemize}
Broj radnika se podrazumevano bira kao $\text{cpu\_count}()-1$, a može se
eksplicitno podesiti preko CLI opcije \texttt{--workers}. Pseudokod paralelne
verzije dat je Algoritmom~\ref{alg:parallel}.

\begin{algorithm}[H]
\caption{GRASP\_Parallel$(D,\, T,\, n_{\text{workers}})$ -- razlika u odnosu na serijsku verziju}
\label{alg:parallel}
\SetAlgoLined
Pokreni \texttt{Pool} od $n_{\text{workers}}$ procesa, pošalji im $D$ jednom\;
$it \leftarrow 0$\;
\While{$it < T$}{
  $\text{batch} \leftarrow \min(n_{\text{workers}},\, T - it)$\;
  Generiši $\text{batch}$ zadataka: konstrukcija$+$VND (svež restart ili ILS),
  a svakih $8$ iteracija i prekrajanje putanja između dva elitna rešenja\;
  \tcp{$\alpha$, ILS semena i roditelji za PR uzimaju se iz stanja na početku batch-a}
  $\text{results} \leftarrow$ paralelno izračunaj sve zadatke iz batch-a\;
  \For{svaki rezultat iz $\text{results}$ (sekvencijalno, istim redosledom kao serijska GRASP petlja)}{
    Ažuriraj skup elitnih rešenja, najbolje rešenje, (svakih 8) prekrajanje
    putanja, (posle 20 zastoja) perturbaciju\;
    $it \mathrel{+}= 1$\;
  }
}
Zatvori \texttt{Pool}\;
\Return najbolje pronađeno rešenje
\end{algorithm}

\subsubsection{Kompletan pseudokod GRASP algoritma}

Kompletan pseudokod (serijski ekvivalent) dat je
Algoritmom~\ref{alg:graspfull}. Paralelna verzija
(Algoritam~\ref{alg:parallel}) izvršava telo petlje po $it$ na isti način po
\textit{redosledu} sekvencijalnog dela, s tim što se konstrukcija$+$VND
(odnosno prekrajanje putanja) za grupu od $n_{\text{workers}}$ uzastopnih
vrednosti $it$ računaju paralelno pre nego što se sekvencijalni deo (ažuriranje
skupa elitnih rešenja, najboljeg rešenja, prekrajanja putanja i perturbacije)
primeni na svaki rezultat iz batch-a, tim redosledom -- pri čemu parametar
$\alpha$, ILS semena i roditelji za prekrajanje putanja odgovaraju stanju na
početku batch-a, a ne stanju koje bi serijska petlja imala nakon prethodnih
iteracija unutar batch-a (videti odeljak~\ref{sec:vektorizacija}).

\begin{algorithm}[H]
\caption{GRASP$(D,\, T)$}
\label{alg:graspfull}
\SetAlgoLined
\KwIn{Podaci instance $D$, broj iteracija $T$}
\KwOut{Najbolje pronađeno rešenje $sol^*$}
$sol^* \leftarrow \varnothing$;\quad
$\text{best\_cost} \leftarrow +\infty$;\quad
$\text{last\_imp} \leftarrow 0$;\quad
$E \leftarrow \varnothing$ (skup elitnih rešenja, veličina 12)\;
\For{$it = 0$ \KwTo $T-1$}{
  \eIf{$it - \text{last\_imp} > 15$}{
    $\alpha \leftarrow \mathcal{U}(0{,}4,\, 0{,}7)$\;
  }{
    $\alpha \leftarrow \mathcal{U}(0{,}05,\, 0{,}5)$\;
  }
  \eIf{$|E| \geq 1$ \textbf{i} nasumično sa verovatnoćom $p_{seed}$}{
    $seed \leftarrow E.\text{random}()$\;
    $mask \leftarrow \text{elite\_dc\_mask}(seed)$ \tcp*{iz neperturbovanog semena}
    $sol \leftarrow$ Perturbacija$(seed, D)$\;
    $sol \leftarrow$ VND$(D,\, sol,\, mask)$\;
  }{
    $mask \leftarrow$ RandomDcMask$(D.J)$\;
    $sol \leftarrow$ GreedyRandomizedConstruction$(D, \alpha, mask)$\;
    $sol \leftarrow$ VND$(D,\, sol,\, mask)$\;
  }
  $E.$add$(sol)$\;
  \If{$sol.\text{cost} < \text{best\_cost}$}{
    $sol^* \leftarrow sol$;\quad
    $\text{best\_cost} \leftarrow sol.\text{cost}$;\quad
    $\text{last\_imp} \leftarrow it$\;
  }
  \If{$it \bmod 8 = 0$ \textbf{i} $|E| \geq 2$}{
    $cand \leftarrow$ PathRelinking$(E.\text{best}(),\, E.\text{random}(),\, D)$\;
    $E.$add$(cand)$\;
    \If{$cand.\text{cost} < \text{best\_cost}$}{
      $sol^* \leftarrow cand$;\quad
      $\text{best\_cost} \leftarrow cand.\text{cost}$;\quad
      $\text{last\_imp} \leftarrow it$\;
    }
  }
  \If{$it - \text{last\_imp} > 20$}{
    $sol' \leftarrow$ RuinAndRecreate$(sol^*, D)$\;
    \If{\textsc{Repair}$(sol', D)$ i $sol'.\text{is\_valid}()$}{
      $sol' \leftarrow$ VND$(D, sol')$\;
      $E.$add$(sol')$\;
      \If{$sol'.\text{cost} < \text{best\_cost}$}{
        $sol^* \leftarrow sol'$;\quad
        $\text{best\_cost} \leftarrow sol'.\text{cost}$;\quad
        $\text{last\_imp} \leftarrow it$\;
      }
    }
    $\text{last\_imp} \leftarrow it$ \tcp*{sprečava uzastopne perturbacije}
  }
}
\Return $sol^*$
\end{algorithm}

% ===============================================================
\section{Rezultati testiranja i analiza}
% ===============================================================

\subsection{Test instance}

Za testiranje je korišćen skup od $60$ instanci podeljenih u tri grupe
(generisane koordinatnim modelom -- čvorovi nasumično u kvadratu
$[-400,400]^2$, varijabilni troškovi su Euklidske distance, fiksni troškovi
su višekratnici tih distanci, $r_f \in [10,15]$ za I. etapu i $r_g \in [5,10]$
za II. etapu):
\begin{itemize}
  \item \textbf{Male} (small, 20 instanci): $I \in [4, 18]$, $J \in [8, 36]$, $K \in [16, 80]$.
  \item \textbf{Srednje} (medium, 20 instanci): $I \in [20, 60]$, $J \in [40, 120]$, $K \in [80, 275]$.
  \item \textbf{Velike} (large, 20 instanci): $I \in [70, 160]$, $J \in [140, 320]$, $K \in [280, 800]$.
\end{itemize}
Svaki dobavljač ima kapacitet $s_i = 80$; potražnja kupaca $d_k \in [10,30]$
je generisana nasumično, a zatim normalizovana tako da ukupna potražnja tačno
odgovara ukupnoj ponudi ($\sum_k d_k = \sum_i s_i$). U trenutku pisanja rada
izvršene su sve male i srednje instance, a od velikih $11$ od $20$
(\texttt{large\_1}--\texttt{large\_11}); preostalih $9$ velikih instanci je u
toku i nije uključeno u proseke i analizu.

\subsection{Referentna rešenja (CPLEX)}

Referentna rešenja dobijena su CPLEX solverom (IBM ILOG CPLEX) korišćenjem MIP
formulacije iz odeljka~1, sa vremenskim limitom od $7200$ s i jednom niti.
Status \texttt{OPTI} označava da je CPLEX dokazao optimalnost; status
\texttt{FEAS} da je dostignut vremenski limit i vraćeno najbolje pronađeno
dopustivo rešenje (bez dokaza optimalnosti) -- to je i dalje referentna
vrednost korišćena za $Bestsol$ (optimalno ili najbolje poznato rešenje). Od
malih instanci $14$ je rešeno optimalno, a $6$ uz status \texttt{FEAS}; sve
srednje i velike instance imaju status \texttt{FEAS}.

\subsection{Podešavanja GRASP-a, platforma i metodologija izveštaja}

\textit{Platforma za GRASP}: x86\_64, 16 CPU jezgara, Python 3.12.3, NumPy
2.4.4. Algoritam je izvršavan paralelno sa
$n_{\text{workers}} = \text{cpu\_count}()-1 = 15$ radnika.

Broj izvršavanja i iteracija zavisi od grupe instanci:
\begin{itemize}
  \item Male instance: $k = 10$ izvršavanja po instanci, $T = 1000$ iteracija;
  \item Srednje instance: $k = 10$ izvršavanja po instanci, $T = 1000$ iteracija;
  \item Velike instance: $k = 3$ izvršavanja po instanci, $T = 100$ iteracija.
\end{itemize}
Označimo sa $D_i$ instancu, a sa $r$-to izvršavanje ($r = 0, 1, \ldots, k-1$)
određeno seed-om $42 + 100r$ (deterministički deo pretrage). Za svako
izvršavanje $i$ algoritam beleži izveštajne veličine:
\begin{itemize}
  \item $sol_i$ -- cena najboljeg pronađenog rešenja;
  \item $t_i$ -- vreme od početka izvršavanja do trenutka prvog nalaska
    $sol_i$ (beleži se u glavnoj petlji, u prekrajanju putanja i u
    perturbaciji, u tačnom trenutku kada se najbolje rešenje poboljša);
  \item $ttot_i$ -- ukupno vreme izvršavanja;
  \item $eval_i$ -- broj poziva funkcije cilja (\texttt{eval\_cost()}, koji se
    poziva i iz \texttt{compute\_cost()} i direktno iz prihvatne probe lokalne
    pretrage), odnosno globalni brojač koji se uvećava pri svakom stvarnom
    izračunavanju vrednosti funkcije cilja; u paralelnom režimu svaki radnik
    ima sopstveni brojač koji se resetuje na početku svakog zadatka i vraća
    kao deo rezultata, a glavni proces ga sabira sa sopstvenim pozivima.
\end{itemize}
$Bestsol$ se određuje kao CPLEX referentna vrednost ako je poznata i bolja od
najboljeg MA rezultata, u suprotnom kao najbolji MA rezultat kroz $k$
izvršavanja; zatim se za svako izvršavanje računa
\[
  gap_i = 100 \cdot \frac{|sol_i - Bestsol|}{|Bestsol|}.
\]
U tabelama su prikazane srednje vrednosti $\bar t$, $\overline{ttot}$,
$\overline{eval}$ i $\overline{gap}$, kao i standardna devijacija $\sigma$ (u
\% gap-a), računate preko $k$ izvršavanja po instanci. Kolona \emph{Best MA} se
razlikuje od ostalih: ona ne prikazuje srednju vrednost, već cenu najboljeg
(najmanjeg) rešenja pronađenog u bilo kom od $k$ izvršavanja za datu instancu.

\subsection{Male instance}

\begin{longtable}{r l r r r r r r r}
\caption{Rezultati za male instance ($T=1000$, $k=10$)} \\
\toprule
\# & $I/J/K$ & Ref.\,(CPLEX) & Best MA & $\bar t$\,(s) & $\overline{ttot}$\,(s) & $\overline{eval}$\,($\times 10^3$) & $\overline{gap}$\,(\%) & $\sigma$\,(\%) \\
\midrule
\endfirsthead
\multicolumn{9}{c}{\tablename\ \thetable{} -- nastavak} \\
\toprule
\# & $I/J/K$ & Ref.\,(CPLEX) & Best MA & $\bar t$\,(s) & $\overline{ttot}$\,(s) & $\overline{eval}$\,($\times 10^3$) & $\overline{gap}$\,(\%) & $\sigma$\,(\%) \\
\midrule
\endhead
\midrule
\multicolumn{9}{r}{\textit{Nastavlja se na sledećoj strani}} \\
\endfoot
\bottomrule
\endlastfoot
1 & 4/8/16 & 150\,568 (OPTI) & 150\,568 & 1.4 & 16.0 & 1270.2 & 0.00 & 0.00 \\
2 & 4/8/20 & 156\,022 (OPTI) & 159\,520 & 3.2 & 18.1 & 1472.7 & 2.24 & 0.00 \\
3 & 4/10/20 & 161\,224 (OPTI) & 163\,375 & 4.3 & 24.3 & 1661.5 & 1.33 & 0.00 \\
4 & 6/12/24 & 186\,537 (OPTI) & 187\,642 & 5.0 & 29.5 & 1832.8 & 1.51 & 0.71 \\
5 & 6/12/30 & 213\,785 (OPTI) & 215\,740 & 20.6 & 39.7 & 2351.2 & 2.13 & 0.61 \\
6 & 6/15/30 & 188\,954 (OPTI) & 194\,359 & 20.7 & 48.0 & 2652.5 & 3.52 & 0.61 \\
7 & 8/16/32 & 231\,984 (OPTI) & 238\,341 & 11.4 & 44.2 & 2478.8 & 3.30 & 0.33 \\
8 & 8/16/40 & 202\,467 (OPTI) & 210\,295 & 43.0 & 84.1 & 3223.0 & 4.61 & 0.50 \\
9 & 8/20/40 & 210\,740 (OPTI) & 214\,410 & 28.4 & 87.7 & 3584.0 & 1.84 & 0.29 \\
10 & 10/20/40 & 258\,236 (OPTI) & 264\,356 & 41.2 & 74.2 & 3457.6 & 3.22 & 0.66 \\
11 & 10/20/50 & 286\,991 (OPTI) & 295\,282 & 51.3 & 107.0 & 4412.0 & 3.48 & 0.44 \\
12 & 10/25/50 & 284\,643 (FEAS) & 297\,645 & 86.3 & 144.9 & 5795.7 & 4.87 & 0.21 \\
13 & 12/24/48 & 278\,044 (OPTI) & 290\,923 & 99.2 & 134.8 & 4898.2 & 4.98 & 0.39 \\
14 & 12/24/60 & 284\,648 (OPTI) & 304\,953 & 79.1 & 141.2 & 6274.8 & 7.66 & 0.39 \\
15 & 12/30/60 & 315\,417 (OPTI) & 328\,859 & 124.0 & 180.4 & 6230.3 & 4.88 & 0.43 \\
16 & 14/28/56 & 349\,820 (FEAS) & 363\,512 & 151.9 & 207.7 & 6465.9 & 4.73 & 0.46 \\
17 & 14/28/70 & 368\,942 (FEAS) & 377\,343 & 127.0 & 244.8 & 8012.0 & 3.17 & 0.49 \\
18 & 16/32/64 & 349\,036 (FEAS) & 362\,127 & 112.7 & 241.4 & 7093.1 & 4.38 & 0.28 \\
19 & 16/32/80 & 351\,345 (FEAS) & 372\,703 & 248.5 & 319.4 & 8372.6 & 6.72 & 0.39 \\
20 & 18/36/72 & 400\,600 (FEAS) & 420\,027 & 216.7 & 306.0 & 7708.9 & 5.97 & 0.55 \\
\midrule
\multicolumn{4}{r}{\textbf{Prosek}} & 73.8 & 124.7 & 4462.4 & \textbf{3.73} & \textbf{0.39} \\
\end{longtable}

\clearpage

\subsection{Srednje instance}

\begin{longtable}{r l r r r r r r r}
\caption{Rezultati za srednje instance ($T=1000$, $k=10$)} \\
\toprule
\# & $I/J/K$ & Ref.\,(CPLEX) & Best MA & $\bar t$\,(s) & $\overline{ttot}$\,(s) & $\overline{eval}$\,($\times 10^3$) & $\overline{gap}$\,(\%) & $\sigma$\,(\%) \\
\midrule
\endfirsthead
\multicolumn{9}{c}{\tablename\ \thetable{} -- nastavak} \\
\toprule
\# & $I/J/K$ & Ref.\,(CPLEX) & Best MA & $\bar t$\,(s) & $\overline{ttot}$\,(s) & $\overline{eval}$\,($\times 10^3$) & $\overline{gap}$\,(\%) & $\sigma$\,(\%) \\
\midrule
\endhead
\midrule
\multicolumn{9}{r}{\textit{Nastavlja se na sledećoj strani}} \\
\endfoot
\bottomrule
\endlastfoot
1 & 20/40/80 & 383\,558 (FEAS) & 398\,577 & 275.8 & 372.2 & 9327.8 & 4.52 & 0.44 \\
2 & 20/40/100 & 415\,983 (FEAS) & 434\,525 & 369.3 & 463.4 & 12731.1 & 5.16 & 0.57 \\
3 & 20/50/100 & 400\,875 (FEAS) & 428\,213 & 347.5 & 565.3 & 13443.9 & 8.17 & 0.65 \\
4 & 25/50/100 & 395\,660 (FEAS) & 414\,693 & 505.3 & 625.2 & 13761.7 & 5.82 & 0.46 \\
5 & 25/50/125 & 474\,978 (FEAS) & 504\,505 & 636.6 & 788.2 & 16568.3 & 7.33 & 0.72 \\
6 & 25/60/120 & 447\,777 (FEAS) & 474\,854 & 890.2 & 969.3 & 19571.9 & 7.11 & 0.76 \\
7 & 30/60/120 & 474\,863 (FEAS) & 498\,115 & 667.6 & 862.1 & 15725.4 & 5.93 & 0.52 \\
8 & 30/60/150 & 514\,195 (FEAS) & 553\,998 & 1796.2 & 2103.5 & 27681.7 & 8.87 & 0.95 \\
9 & 30/75/150 & 503\,494 (FEAS) & 530\,874 & 1520.2 & 1800.4 & 27182.2 & 6.51 & 0.59 \\
10 & 35/70/140 & 525\,261 (FEAS) & 563\,494 & 2024.9 & 2139.8 & 25283.5 & 7.78 & 0.30 \\
11 & 35/70/175 & 551\,427 (FEAS) & 576\,803 & 1984.1 & 2144.5 & 30091.8 & 5.98 & 0.88 \\
12 & 40/80/160 & 478\,570 (FEAS) & 503\,072 & 2169.6 & 2519.7 & 26760.4 & 6.25 & 0.77 \\
13 & 40/80/200 & 616\,913 (FEAS) & 669\,048 & 3196.6 & 3527.8 & 41645.2 & 9.73 & 0.76 \\
14 & 45/90/180 & 689\,812 (FEAS) & 741\,059 & 3013.7 & 3846.8 & 40371.7 & 8.13 & 0.56 \\
15 & 45/90/225 & 557\,928 (FEAS) & 604\,169 & 3068.2 & 3942.8 & 42065.7 & 8.91 & 0.44 \\
16 & 50/100/200 & 592\,227 (FEAS) & 648\,696 & 3706.4 & 4128.7 & 39819.4 & 10.64 & 0.64 \\
17 & 50/100/250 & 612\,725 (FEAS) & 652\,416 & 5573.9 & 6418.6 & 59909.0 & 7.71 & 0.65 \\
18 & 55/110/220 & 685\,116 (FEAS) & 753\,348 & 5187.0 & 5969.0 & 52562.4 & 11.27 & 0.67 \\
19 & 55/110/275 & 580\,226 (FEAS) & 626\,831 & 6643.6 & 7732.8 & 64967.1 & 8.93 & 0.58 \\
20 & 60/120/240 & 677\,683 (FEAS) & 729\,939 & 6121.0 & 6861.8 & 52925.8 & 8.59 & 0.44 \\
\midrule
\multicolumn{4}{r}{\textbf{Prosek}} & 2484.9 & 2889.1 & 31619.8 & \textbf{7.67} & \textbf{0.62} \\
\end{longtable}

\clearpage

\subsection{Velike instance}

Velike instance predstavljaju najzahtevniju grupu; u trenutku pisanja rada
izvršeno je $11$ od $20$ instanci (\texttt{large\_1}--\texttt{large\_11}) sa
$k=3$ izvršavanja po instanci i $T=100$ iteracija. Preostalih $9$ instanci je u
toku i nije uključeno u proseke i analizu.

\begin{longtable}{r l r r r r r r r}
\caption{Rezultati za velike instance ($T=100$, $k=3$)} \\
\toprule
\# & $I/J/K$ & Ref.\,(CPLEX) & Best MA & $\bar t$\,(s) & $\overline{ttot}$\,(s) & $\overline{eval}$\,($\times 10^3$) & $\overline{gap}$\,(\%) & $\sigma$\,(\%) \\
\midrule
\endfirsthead
\multicolumn{9}{c}{\tablename\ \thetable{} -- nastavak} \\
\toprule
\# & $I/J/K$ & Ref.\,(CPLEX) & Best MA & $\bar t$\,(s) & $\overline{ttot}$\,(s) & $\overline{eval}$\,($\times 10^3$) & $\overline{gap}$\,(\%) & $\sigma$\,(\%) \\
\midrule
\endhead
\midrule
\multicolumn{9}{r}{\textit{Nastavlja se na sledećoj strani}} \\
\endfoot
\bottomrule
\endlastfoot
1 & 70/140/280 & 741\,988 (FEAS) & 829\,048 & 1181.7 & 1220.4 & 7665.7 & 12.98 & 0.92 \\
2 & 70/140/350 & 825\,845 (FEAS) & 954\,277 & 1623.8 & 1901.5 & 10823.2 & 16.10 & 0.43 \\
3 & 80/160/320 & 851\,981 (FEAS) & 999\,912 & 1781.8 & 2065.9 & 10358.8 & 17.97 & 0.45 \\
4 & 80/160/400 & 807\,109 (FEAS) & 909\,778 & 2765.4 & 2889.6 & 13246.1 & 12.87 & 0.17 \\
5 & 90/180/360 & 801\,015 (FEAS) & 902\,199 & 2408.0 & 2650.7 & 11503.0 & 13.37 & 0.55 \\
6 & 90/180/450 & 953\,624 (FEAS) & 1\,075\,114 & 3242.3 & 4157.6 & 16344.4 & 13.83 & 0.78 \\
7 & 100/200/400 & 1\,081\,573 (FEAS) & 1\,237\,060 & 4196.2 & 4437.0 & 16028.4 & 15.32 & 1.06 \\
8 & 100/200/500 & 916\,397 (FEAS) & 1\,066\,027 & 4872.1 & 6323.7 & 19903.8 & 17.63 & 0.92 \\
9 & 110/220/440 & 871\,261 (FEAS) & 1\,022\,419 & 5708.3 & 5708.3 & 16513.4 & 18.36 & 1.00 \\
10 & 110/220/550 & 945\,683 (FEAS) & 1\,080\,777 & 7082.7 & 8055.4 & 21771.7 & 15.29 & 0.93 \\
11 & 120/240/480 & 971\,579 (FEAS) & 1\,112\,135 & 8410.9 & 8410.9 & 21585.8 & 15.55 & 0.86 \\
\midrule
\multicolumn{4}{r}{\textbf{Prosek}} & 3933.9 & 4347.4 & 15067.6 & \textbf{15.39} & \textbf{0.73} \\
\end{longtable}

\subsection{Analiza rezultata}

\subsubsection{Analiza -- male instance}

Na malim instancama GRASP postiže prosečno odstupanje od $\overline{gap} =
3{,}73\%$ ($\sigma = 0{,}39\%$) u odnosu na CPLEX referentnu vrednost, uz
$k=10$ nezavisnih izvršavanja po instanci. Na instanci~1 ($4 \times 8 \times
16$) algoritam u svih 10 izvršavanja pronalazi optimalno rešenje (gap
$0{,}00\%$), a na instancama~2 i~3 gap je manji od $2{,}5\%$. Standardna
devijacija je mala (uglavnom ispod $0{,}7\%$), što ukazuje na konzistentno
ponašanje algoritma za fiksirani budžet od $T=1000$ iteracija. Najveći gap
postiže se na instanci~14 ($12 \times 24 \times 60$, $7{,}66\%$), što ukazuje
da porast broja kupaca po DC-u otežava lokalnoj pretrazi da u zadatom broju
iteracija pronađe globalni optimum. Prosečno vreme do najboljeg rešenja
$\bar t$ raste sa dimenzijama instance (od $1{,}4$~s do $248{,}5$~s), dok je
ukupno vreme $\overline{ttot}$ svugde ispod $320$~s.

\subsubsection{Analiza -- srednje instance}

Na srednjim instancama prosečno odstupanje iznosi $\overline{gap} = 7{,}67\%$
($\sigma = 0{,}62\%$). Sve CPLEX referentne vrednosti za srednje instance
imaju status \texttt{FEAS} -- CPLEX je dostigao vremenski limit od $7200$~s bez
dokaza optimalnosti, pa referentna vrednost nije nužno optimalna, već najbolje
pronađeno dopustivo rešenje kroz snažnu $B\&B$/\textit{cutting-plane} pretragu
-- suštinski drugačiji (i po vremenu mnogo skuplji) mehanizam od GRASP lokalne
pretrage sa $T=1000$ iteracija. Najmanji gap je na instanci~1
($20\times40\times80$, $4{,}52\%$), a najveći na instanci~18
($55\times110\times220$, $11{,}27\%$). Vreme izvršavanja raste sa dimenzijama:
$\overline{ttot}$ ide od $372$~s (instanca~1) do $7733$~s (instanca~19, preko
2 sata po izvršavanju), a broj evaluacija funkcije cilja $\overline{eval}$
raste do približno $65$ miliona poziva po izvršavanju.

\subsubsection{Analiza -- velike instance}

Na velikim instancama (11 izvršenih) prosečno odstupanje iznosi
$\overline{gap} = 15{,}39\%$ ($\sigma = 0{,}73\%$), nastavljajući trend rasta
uočen kod malih i srednjih grupa ($3{,}73\% \to 7{,}67\% \to 15{,}39\%$) --
očekivano, budući da je broj iteracija ($T=100$) najmanji od sve tri grupe,
dok je prostor pretrage ($I,J,K$ do $120,240,480$ među izvršenim instancama)
najveći. Najmanji gap je na instanci~4 ($80\times160\times400$, $12{,}87\%$),
a najveći na instanci~9 ($110\times220\times440$, $18{,}36\%$); razlike između
instanci ukazuju da na ovim dimenzijama varijabilnost zavisi od konkretne
strukture instance, ne samo od njene veličine. Vreme izvršavanja raste
dramatično: $\overline{ttot}$ ide od $1220$~s (instanca~1) do preko $8400$~s
(instanca~11, oko 2,5 sata po izvršavanju). Očekuje se da preostalih 9
instanci (do $160 \times 320 \times 800$) bude znatno zahtevnije.

\subsubsection{Zajednički zaključci}

\begin{itemize}
  \item Odstupanje GRASP-a od CPLEX referentne vrednosti raste sa dimenzijom
    instance ($3{,}73\% \to 7{,}67\% \to 15{,}39\%$), što je očekivana
    posledica fiksnog broja iteracija po grupi ($1000 \to 1000 \to 100$) uz
    brzo rastući prostor pretrage.
  \item Standardna devijacija $\sigma$ je kroz sve tri grupe relativno mala
    (prosečno $0{,}39$--$0{,}73\%$, pojedinačno najviše do $\approx 1{,}1\%$),
    što pokazuje da je algoritam, uz višestruka izvršavanja po instanci,
    dovoljno konzistentan -- razlike između pojedinačnih izvršavanja iste
    instance znatno su manje od razlika između instanci.
  \item Za srednje i velike instance (i deo malih) CPLEX referentne vrednosti
    same po sebi nisu dokazano optimalne (status \texttt{FEAS}), pa je stvarni
    gap GRASP-a prema (nepoznatom) optimumu verovatno nešto manji od
    prikazanog.
  \item Vektorizacija i paralelizacija (odeljci 2.2.9--2.2.10) preduslov su da
    je $k=10$ izvršavanja po svih malih i srednjih instanci uopšte bilo
    izvodljivo u razumnom vremenu; egzaktno filtriranje kandidata dodatno
    ubrzava lokalnu pretragu bez pogoršanja kvaliteta.
\end{itemize}

% ===============================================================
\section{Zaključak}
% ===============================================================

U ovom radu predložena je GRASP metaheuristika za Dvostepeni transportni
problem sa fiksnim troškovima (TS-FCTP). Algoritam kombinuje pohlepno-nasumičnu
konstrukciju sa adaptivno podešavanim parametrom $\alpha$, VND lokalnu pretragu
sa pet okolina, egzaktno filtriranje kandidata po zatvorenoj formuli za trošak
druge etape, skup elitnih rešenja sa raznolikošću po Jaccardovom rastojanju,
prekrajanje putanja razmenom kolona kupaca (u paralelnoj verziji izvršavan u
radnicima), ILS perturbaciju elitnih rešenja ograničenu maskom na DC
konfiguraciju semena i ruin-and-recreate perturbaciju pri dugotrajnom zastoju.

Centralna arhitekturalna odluka -- direktna pretraga u prostoru $y$-tokova uz
rekonstrukciju $x$-tokova procedurom \textsc{Repair} -- efikasno smanjuje
dimenzionalnost problema. Nad tom osnovom, implementacija je dodatno
vektorizovana (NumPy) i paralelizovana (\texttt{multiprocessing}, batch
nezavisnih iteracija), što je omogućilo višestruka izvršavanja po instanci i
dobijanje statistički smislenih ocena ($\overline{gap}$, $\sigma$).

Testiranje pokazuje da GRASP postiže dobra rešenja na malim instancama
(prosečan gap $3{,}73\%$, $\sigma = 0{,}39\%$, sa tačno optimalnim rešenjem na
najmanjoj instanci), uz umeren pad kvaliteta na srednjim instancama (gap
$7{,}67\%$, $\sigma = 0{,}62\%$) i dalji pad na velikim instancama (gap
$15{,}39\%$, $\sigma = 0{,}73\%$, na 11 od 20 izvršenih instanci), usled
fiksnog i sve manjeg broja iteracija u odnosu na sve veći prostor pretrage na
krupnijim problemima. Kroz sve tri grupe standardna devijacija ostaje relativno
mala, što ukazuje na konzistentno ponašanje algoritma nezavisno od dimenzije
instance.

Mogući pravci daljeg razvoja:
\begin{itemize}
  \item dovršiti izvršavanje preostalih $9$ velikih instanci i uporediti
    rezultate na kompletnom skupu;
  \item povećati broj iteracija $T$ na velikim instancama (npr.\ dodatnim
    ubrzanjem \textsc{Repair}-a) kako bi se smanjio uočeni gap na toj grupi;
  \item istražiti varijante prekrajanja putanja (npr.\ duži postupak, izbor koraka
    prema uticaju na I. etapu) kako bi se međutačke bolje iskoristile na
    velikim instancama.
\end{itemize}

% ===============================================================
\begin{thebibliography}{9}

\bibitem{jb2009}
N.~Jawahar, A.~N.~Balaji,
\textit{A genetic algorithm for the two-stage supply chain distribution
problem associated with a fixed charge},
European Journal of Operational Research \textbf{194}(2), 496--537 (2009).

\bibitem{calvete2018}
H.~I.~Calvete, C.~Gal\'e, J.~A.~Iranzo, P.~Toth,
\textit{A matheuristic for the two-stage fixed-charge transportation problem},
Computers \& Operations Research \textbf{95}, 113--122 (2018).

\bibitem{balaji2010}
A.~N.~Balaji, N.~Jawahar,
\textit{A simulated annealing algorithm for a two-stage fixed charge
distribution problem of a supply chain},
International Journal of Operational Research \textbf{7}(2), 192--215 (2010).

\bibitem{raj2012}
K.~A.~A.~D.~Raj, C.~Rajendran,
\textit{A genetic algorithm for solving the fixed-charge transportation model:
two-stage problem},
Computers \& Operations Research \textbf{39}(9), 2016--2032 (2012).

\bibitem{pop2016}
P.~C.~Pop, O.~Matei, C.~Pop~Sitar, I.~Zelina,
\textit{A hybrid based genetic algorithm for solving a capacitated fixed-charge
transportation problem},
Carpathian Journal of Mathematics \textbf{32}(2), 225--232 (2016).

\bibitem{feo1995}
T.~A.~Feo, M.~G.~C.~Resende,
\textit{Greedy randomized adaptive search procedures},
Journal of Global Optimization \textbf{6}(2), 109--133 (1995).

\bibitem{eksioglu2003}
B.~Eksioglu, A.~Migdalas, P.~M.~Pardalos,
\textit{Heuristic approaches to production-inventory-distribution problems in
supply chains},
in: Optimization and Industry: New Frontiers, Kluwer Academic Publishers,
15--37 (2003).

\end{thebibliography}

\end{document}
